% =====================================================================
%  Scale-Aware Joint Compression -- Abstract, Introduction, Related Work
%
%  Required preamble packages:
%    \usepackage{booktabs, amsmath, amssymb, algorithm, algpseudocode}
%    \usepackage{multirow, siunitx}
%
%  Section order in the assembled manuscript:
%    frontmatter.tex  (abstract, S1 introduction, S2 related work)
%    methodology.tex  (S3 methodology)
%    results.tex      (S4 results, S5 discussion, S6 limitations, S7 conclusion)
% =====================================================================

\begin{abstract}
Pruning and quantisation are conventionally applied in sequence: prune, recover, then
quantise. Joint formulations, in which the pruning decision is made with knowledge of the
quantisation grid, are reported to improve on this --- but almost always at a single model
size, leaving unanswered the question a practitioner actually faces: at \emph{my} model
scale, is the additional pipeline complexity justified? We report a pre-registered,
confirmatory measurement of \emph{joint gain} --- the difference in perplexity retention
between joint and sequential pipelines at a matched compression budget --- across the Pythia
suite from 160M to 1B targeted parameters, at two precisions, over eight paired calibration
replicates, with optimisation budgets matched by construction and asserted in code. On a
held-out test split accessed once, \textbf{no cell meets our pre-registered
practical-importance criterion} of $\geq\!\SI{1.0}{pp}$ mean gain with a consistent sign.
The largest effects are $+\SI{1.01}{pp}$ at 160M ($7/8$ replicates positive) and
$+\SI{0.93}{pp}$ at 410M ($8/8$, Holm-adjusted $p = 0.047$), each failing a different
criterion by a narrow margin; at 1B the gain falls to $+\SI{0.13}{pp}$. The motivating
hypothesis --- that joint optimisation pays off more at scale --- is not supported. Layerwise
diagnostics reveal why the aggregate picture is not simply a weakening mechanism: mask
divergence and local layer-objective advantage are \emph{undiminished} at 1B, yet the
end-to-end gain collapses, a dissociation we characterise but do not close. We further find
that the precision-specific structure of the effect replicates in a second model family
(Qwen2.5-0.5B), that $30\%$ unstructured sparsity yields no commensurate CPU speedup, and
that a short recovery phase applied equally to both arms \emph{reverses} the ordering in
$8/8$ replicates. We release code, configurations and all run records.
\end{abstract}

\section{Introduction}
\label{sec:intro}

Language models in the $10^{8}$--$10^{9}$ parameter range are increasingly deployed on
commodity CPUs, where compression is not an optimisation but a precondition for running the
model at all. The two dominant techniques --- unstructured pruning and post-training weight
quantisation --- are complementary, and in practice are applied \emph{sequentially}: the
model is pruned, its quality partially recovered, and the result quantised. This ordering is
an engineering convention rather than a derived optimum, and it has a clear structural
deficiency. When the mask is chosen first, it is chosen in ignorance of the quantisation
grid: a weight retained for its magnitude may be rounded to a grid point that renders it
nearly useless, while a weight discarded as small may have survived quantisation intact.

\emph{Joint} formulations address this by making the mask decision under quantised weights
and re-estimating the quantisation scales after the mask moves, so that the two constraints
inform one another within a single reconstruction objective. Such methods are more expensive
to implement, more fragile to tune, and --- critically for a practitioner deciding whether to
adopt one --- are almost universally evaluated at a single model size.

This leaves a gap that matters. Compression sensitivity is known to vary with scale: larger
models carry more redundancy, and the same recipe that devastates a 160M model may barely
perturb a 1B one. If joint optimisation confers a growing advantage as models scale, the
engineering investment is easy to justify; if the advantage shrinks, the sequential pipeline
is the rational default precisely where compression matters most. Neither claim has been
tested under controlled conditions.

A second gap runs alongside it. Compression papers report theoretical parameter reduction and
bit-width savings far more often than measured latency on the hardware the model will
actually run on. A $30\%$ sparse, 4-bit model implies a large nominal compression ratio; what
it delivers on a dense CPU GEMM kernel is a separate empirical question, and one that a FLOP
count cannot answer.

\paragraph{Contributions.}
\begin{enumerate}
  \item \textbf{A controlled measurement of joint gain across an order of magnitude of
        model scale} (160M $\rightarrow$ 410M $\rightarrow$ 1B targeted parameters), with
        dataset, tokeniser, calibration draw, module coverage and optimisation budget held
        fixed by construction rather than by convention.
  \item \textbf{Fairness enforced in code.} Both arms are call-orderings over one shared
        layerwise solver; equality of solver budget is asserted from the configuration
        before a sweep begins and from the persisted records afterwards. We report the
        specific failure modes these assertions were introduced to catch.
  \item \textbf{A pre-registered confirmatory design.} Budgets, sequential orderings and all
        exploratory estimates live on a validation split; the test split was accessed once,
        under a frozen configuration, with a practical-importance threshold fixed in advance.
        We report that \emph{no cell meets it}, including two near-misses that we decline to
        accommodate by moving the bar.
  \item \textbf{A mechanistic dissociation.} Layerwise diagnostics show the joint mechanism
        is \emph{not} weaker at 1B --- mask divergence and local objective advantage are at
        or above their 160M values --- while end-to-end gain collapses by roughly $8\times$.
        We identify two candidate explanations, including a depth confound intrinsic to the
        Pythia suite, and test neither, stating the dissociation as open.
  \item \textbf{Measured rather than inferred deployment claims}, on CPU, with prefill and
        decode separated, and an explicit account of which arms cannot be timed at all and
        why.
  \item \textbf{Three replication checks that a positive result would not have survived:} an
        external-validity leg in a second model family, a matched-quality mechanistic
        control, and a post-hoc recovery ablation that reverses the ordering.
\end{enumerate}

\paragraph{Summary of findings.} On the confirmatory split, joint compression yields a small
positive gain over the frozen sequential baseline at 4-bit precision --- $+\SI{1.01}{pp}$ at
160M and $+\SI{0.93}{pp}$ at 410M --- and effectively nothing at 1B ($+\SI{0.13}{pp}$). At
8-bit precision, where the mechanism is predicted to be inert, gain is indistinguishable from
zero at the two smaller scales and reliably \emph{negative} at 1B. No cell satisfies the
pre-registered practical-importance criterion. The observed direction across scale is
opposite to the motivating hypothesis, but the cross-scale decline is not itself
statistically established, and is confounded with a non-monotone depth profile in the model
suite. We therefore report a negative result on practical importance, a robust
precision-specific structural finding that replicates across model families, and an open
mechanistic question.

\section{Background and related work}
\label{sec:related}

\paragraph{Unstructured pruning and post-training reconstruction.} Magnitude pruning removes
weights of smallest absolute value and remains a strong baseline. Activation-aware criteria
improve on it by weighting magnitude with the energy of the corresponding input channel,
recognising that a large weight multiplying a near-zero activation contributes little. A
parallel line of work reconstructs the surviving weights to preserve layer \emph{outputs}
rather than weight values, solving a least-squares problem in the input Gram matrix
$H = X^{\top}X$; this error-compensation step accounts for most of the benefit over naive
masking, and second-order variants extend it with Hessian-informed column sweeps. Our
solver sits in this family. We deliberately adopt a first-order saliency criterion and a
damped alternating refinement solver rather than a state-of-the-art pruning criterion: a
stronger base method would raise \emph{both} arms, and the quantity of interest here is the
difference between pipelines, not their absolute quality.

\paragraph{Post-training quantisation.} Weight-only post-training quantisation maps weights
to a low-precision grid using scales estimated from calibration statistics, with symmetric
per-channel or per-group schemes trading calibration cost against fidelity. Below 8 bits,
reconstruction against calibration activations becomes necessary rather than optional.
Quantisation-aware training, which propagates gradients through a straight-through estimator,
generally recovers more quality but requires a full training loop --- a memory-compute
bottleneck that rules it out at $10^{9}$ parameters on a single commodity accelerator, and
the reason this study adopts layerwise post-training reconstruction instead.

\paragraph{Combining pruning and quantisation.} The sequential composition of the two is
standard practice. Joint formulations --- variously described as pruning-aware quantisation,
sparse-quantised optimisation, or unified compression --- report advantages from letting the
mask decision see the grid. Our reading of this literature is that the reported gains are
credible but narrowly scoped: they are typically demonstrated at one model size, often
without matched optimisation budgets between the compared pipelines, and rarely with
replicate-level variance. The last point is not a minor omission. As we show in
\S\ref{sec:results}, replicate-to-replicate spread at 4-bit precision is comparable to the
effect itself, so a single-draw comparison cannot distinguish a genuine advantage from a
favourable calibration draw.

\paragraph{Scale as an experimental variable.} The Pythia suite is designed for controlled
scale studies: models across the size range share data ordering, tokeniser and training
recipe, so scale can be varied while holding the training pipeline fixed. This makes it the
natural substrate for our question. It is not, however, a perfectly isolated scale axis ---
depth, width, head configuration and tokens-per-parameter all covary along the suite, and we
document in \S\ref{subsec:limits} a specific consequence: \emph{pythia-1b is shallower than
pythia-410m}, which confounds the one large movement in our trend.

\paragraph{Sparsity and realised speedup.} A recurring gap in the compression literature is
between nominal and realised benefit. Unstructured sparsity does not accelerate a dense GEMM
kernel: the multiply-accumulates still execute, they merely multiply by zero. Realising a
speedup requires either a sparse kernel or a structured pattern ($N{:}M$) the hardware can
exploit, and CPU support for semi-structured sparsity is considerably less mature than its
GPU counterpart. We therefore treat ``does sparsity become speed?'' as an empirical question
to be measured on the deployment target, and report the answer we obtained rather than the
one the compression ratio implies.

\paragraph{Positioning.} We introduce no new compression method. Our contribution is a
measurement, under controls the existing literature does not consistently apply, of a design
choice that is widely made and rarely tested across scale.
